\chapter[Algebraic States and Formal Correlations]{Algebraic States, Correlations, and Formal Vertices}
\label{app:states_correlations}

\section{Status and scope}

This appendix records the proposed uncompressed oscillator and charge-algebra
constructions used by Chapter~\ref{ch:correlations}. Some related declarations
are present in the CH19 staging files, and the CH19 proof batch is recorded as
Phase-B complete; neither fact is Core promotion. The constructions and
derivations below are therefore mathematical manuscript material, not a claim
that the current Lean Core certifies an abstract state, every displayed
correlator, or a comparison with finite CAR states. The distinction between
the two carriers is part of the result's meaning.

\section{Polynomial oscillator carrier and its vacuum functional}

Fix a finite set of positive modes $1\leq m\leq M$ and two branches
$\nu\in\{R,L\}$. Let $\mathcal P=\mathbb C[X_{\nu m}]$ be the polynomial
algebra in these variables. Creation acts by multiplication,
$C_{\nu m}p=X_{\nu m}p$, while annihilation acts by
$A_{\nu m}p=m\,\partial p/\partial X_{\nu m}$. On monomials define
\begin{equation}
 \langle X^\alpha,X^\beta\rangle
 =\delta_{\alpha,\beta}\prod_{\nu,m}\alpha_{\nu m}!
                                      m^{\alpha_{\nu m}}.
 \label{eq:a08-polynomial-gram}
\end{equation}
Every diagonal weight is positive, so this is a positive definite Hermitian
form on the polynomial carrier. Directly comparing the coefficient of each
monomial shows $A_{\nu m}^*=C_{\nu m}$ on polynomials. Differentiating a
monomial also gives
\begin{equation}
 [A_{\nu m},C_{\eta n}]
   =\delta_{\nu\eta}\delta_{mn}mI,
 \qquad [A_{\nu m},A_{\eta n}]=[C_{\nu m},C_{\eta n}]=0.
 \label{eq:a08-polynomial-ccr}
\end{equation}
Thus finitely many modes generate an infinite-dimensional polynomial space;
there is no energy compression in this construction.

Let $\mathbf 1$ denote the constant polynomial. For a represented word $a$
define $\omega_0(a)=\langle\mathbf 1,a\mathbf 1\rangle$. The vacuum has unit
norm. For any polynomial operator $T$ in the represented generators,
\begin{equation}
 \omega_0(T^*T)=\langle T\mathbf 1,T\mathbf 1\rangle\geq0.
 \label{eq:a08-vacuum-positive}
\end{equation}
The same calculation shows why this is a positive functional on the represented
algebra. To define it on a quotient by CCR relations, one must additionally
show that the representation annihilates the relation ideal; this is a
well-definedness obligation, not a consequence of writing the CCR symbols.

\section{Dressing and the abstract state}

For real $c,s$ with $c^2-s^2=1$, define a modewise transformation by
\begin{equation}
 \beta(C_R)=cC_R+sA_L,\qquad
 \beta(C_L)=cC_L+sA_R,\qquad
 \beta(A_\nu)=\beta(C_\nu)^*.
 \label{eq:a08-beta}
\end{equation}
The inverse is obtained by replacing $s$ by $-s$. If this map preserves the
algebra relations and star operation, the dressed functional is
\begin{equation}
 \omega_P(a)=\omega_0(\beta^{-1}(a)).
 \label{eq:a08-dressed-state}
\end{equation}
It is normalized and positive because $\beta^{-1}$ is a unital star
automorphism: $\omega_P(a^*a)=\omega_0(\beta^{-1}(a)^*\beta^{-1}(a))\geq0$.
The inverse in~\eqref{eq:a08-dressed-state} fixes the state convention; changing
the pullback direction changes which dressed operators annihilate the vacuum.

\section{Both ordered contractions}

Writing $C_R,A_R$ for one fixed mode and suppressing its index, the inverse
transform is $\beta^{-1}(C_R)=cC_R-sA_L$ and
$\beta^{-1}(A_R)=cA_R-sC_L$. In the product vacuum, different branches
factor and $A\mathbf1=0$. Using $A_RC_R\mathbf1=m\mathbf1$ gives
\begin{equation}
 \omega_P(C_RA_R)=ms^2,\qquad
 \omega_P(A_RC_R)=mc^2.
 \label{eq:a08-bare-contractions}
\end{equation}
The second ordered contraction retains its $mc^2$ term even at $s=0$; it is
the CCR contribution. For
\begin{equation}
 \rho_R(x)=L^{-1}\sum_{m=1}^M
   \left(\zeta^{-mx}C_{Rm}+\zeta^{mx}A_{Rm}\right),
 \qquad \zeta=e^{2\pi i/L},
\end{equation}
expansion of the ordered product and mode orthogonality then yields
\begin{equation}
 \omega_P(\rho_R(x)\rho_R(y))
 =L^{-2}\sum_{m=1}^M m\left(
   s^2\zeta^{-m(x-y)}+c^2\zeta^{m(x-y)}\right).
 \label{eq:a08-density-kernel}
\end{equation}
This kernel is generally complex. Replacing it by a cosine would describe a
symmetrized expression and discard the ordering information.

\section{Formal vertices and coefficientwise evaluation}

Let $\mathcal A[[t]]$ be the formal power series algebra over the oscillator
and charge algebra, with central $t=t^*$. For $X\in\mathcal A$, define
\begin{equation}
 \exp_{\mathrm{form}}(itX)=\sum_{n\geq0}\frac{(itX)^n}{n!}.
 \label{eq:a08-formal-exp}
\end{equation}
Each coefficient is a finite algebraic expression. Extend a state coefficient
by coefficient to a map into $\mathbb C[[t]]$; this is a formal evaluator,
not an ordinary complex-valued positive state on all series, and it does not
permit setting $t=1$ without convergence information.

With opposite spatial orientations for the two branches, let $Q_R(x),Q_L(x)$
be the corresponding finite Hermitian current sums, and put
\begin{equation}
 X_{\mathrm{CDW}}=(c-s)(Q_R+Q_L),\qquad
 X_{\mathrm{SC}}=(c+s)(Q_R-Q_L).
\end{equation}
Let $K_J(x)$ be the specified charge-shift and source-phase factor for
$J\in\{\mathrm{CDW},\mathrm{SC}\}$. The Weyl-type formal vertex is
\begin{equation}
 W_J(P,t,x)=L^{-1}K_J(x)
                 \exp_{\mathrm{form}}(itX_J(P,x)).
 \label{eq:a08-vertex}
\end{equation}
The order of $K_J$ and the exponential is part of the definition. These
vertices are not the finite projected Mattis--Mandelstam maps of
Chapter~\ref{ch:mandelstam}, nor are they silently normal-ordered
exponentials.

For $D_1(x,y)=\sum_{m=1}^M[1-\cos(2\pi m(x-y)/L)]/m$, the opposite
orientations cancel the two chiral scalar commutators. The polynomial vacuum
moment calculation gives variances $4gD_1$ and $4g^{-1}D_1$ for the respective
field differences. To see how the Gaussian factor follows, write
$Y_J=X_J(y)-X_J(x)$ and $V_J=\omega_{P,N}(Y_J^2)$. The vacuum is annihilated
by each $A_{\nu m}$, so in expanding a word in $Y_J$ one moves annihilators
to the right; each commutation contributes its scalar CCR term. Pairing the
leftmost factor with each of the remaining $2n-1$ factors gives
\[
 \omega_{P,N}(Y_J^{2n})=(2n-1)!!\,V_J^n,\qquad
 \omega_{P,N}(Y_J^{2n+1})=0.
\]
Therefore
\[
 \omega_{P,N}(\exp_{\mathrm{form}}(itY_J))
 =\sum_{n\geq0}\frac{(-1)^n V_J^n t^{2n}}{2^n n!}
 =\operatorname{Exp}_{\mathrm{form}}(-V_Jt^2/2).
\]
Substituting $V_{\mathrm{CDW}}=4gD_1$ and
$V_{\mathrm{SC}}=4g^{-1}D_1$, and multiplying the ordered charge phases,
gives the proposed coefficientwise identities
\begin{align}
 \omega_{P,N}(W_{\mathrm{CDW}}(t,x)^*W_{\mathrm{CDW}}(t,y))
 &=L^{-2}\zeta^{(N_L-N_R-1)(y-x)}
       \operatorname{Exp}_{\mathrm{form}}(-2gD_1(x,y)t^2),\\
 \omega_{P,N}(W_{\mathrm{SC}}(t,x)^*W_{\mathrm{SC}}(t,y))
 &=L^{-2}\zeta^{(N_R+N_L)(y-x)}
       \operatorname{Exp}_{\mathrm{form}}(-2g^{-1}D_1(x,y)t^2).
 \label{eq:a08-formal-correlators}
\end{align}
This calculation is coefficientwise: for each power of $t$, only finitely
many contractions enter. It establishes formal identities only: no analytic evaluation
at $t=1$, continuum power law, or finite-CAR comparison follows.

\section{Why compression does not inherit the abstract vacuum}

Compressing operators to a finite-dimensional subspace changes their products:
in general $PAP\,PCP\ne PACP$. In particular, a dressed annihilator assembled
from compressed creation and annihilation operators need not have a common
kernel. A useful finite warning is the two-mode total-degree-at-most-one slice
with basis $(1,X_R,X_L)$. For $c=5/3$, $s=4/3$, the compressed combinations
$Q_R=cD_R+s\widetilde X_L$ and $Q_L=cD_L+s\widetilde X_R$ give
\begin{equation}
 Q_R^*Q_R+Q_L^*Q_L=\operatorname{diag}(32/9,25/9,25/9).
 \label{eq:a08-compressed-matrix}
\end{equation}
The diagonal entries are positive, so the common kernel is zero: if both
$Q_Rv=0$ and $Q_Lv=0$, taking the inner product of~\eqref{eq:a08-compressed-matrix}
with $v$ forces $v=0$. This rules out a normalized vacuum satisfying those
two compressed annihilation equations on this slice. It does not rule out a
ground state for a separately defined finite Hamiltonian.

For clarity, the matrix statement can be checked directly. On the ordered basis
$(1,X_R,X_L)$, compression of multiplication discards degree-two terms, so
$Q_R1=sX_L$, $Q_RX_R=c$, $Q_RX_L=0$; similarly
$Q_L1=sX_R$, $Q_LX_R=0$, and $Q_LX_L=c$. Thus the squared norms of the three
columns of $Q_R$ are $(s^2,c^2,0)$ and those of $Q_L$ are $(s^2,0,c^2)$.
Their sum is diagonal with entries $(2s^2,c^2,c^2)=(32/9,25/9,25/9)$ for the
chosen parameters. Positivity of all three entries proves the common-kernel
claim without a spectral theorem.

\begin{remark}[Formalization boundary]
The polynomial model, charge lattice, state pullback, coefficientwise
correlators, and the compressed finite example have distinct carriers and
proof obligations. The CH19 staging status does not establish that every
construction here has passed Core promotion. Power-law asymptotics, analytic
exponentials, and equality with finite fermionic expectations remain separate
problems.
\end{remark}
